\chapter{本書の構成と実行例の読み方}
\label{chap:usage}

\section{対象と到達目標}

本書の対象は、Python を初めて学ぶ読者に加え、短い解析コードは書けても、実行環境の管理、結果の検証、再利用可能な構成、高速化に不安を持つ読者である。到達目標は、読者が次の要素を一つの解析手順として組み立て、各選択の理由を説明できるようになることである。

\begin{itemize}
\item シェルによるファイル・ディレクトリの操作と、通常の実行結果（標準出力）・エラーメッセージ（標準エラー出力）の保存
\item Python のバージョン、仮想環境、依存ライブラリの管理と記録
\item テキスト、表、バイナリデータの読み込みと、型、欠損値、単位、時刻系の確認
\item 配列・表の整形と集計、および図、統計量、フィット結果に基づくデータの評価
\item テストと測定に基づくコードの整理と、解析結果の同一性を保った高速化
\end{itemize}

既習事項が多い読者は、必要な章から読み始めてもよい。ただし、後半の例は、それ以前の章で導入した用語とファイルを前提としている。分からない用語は索引から説明箇所を探せる。索引の太字のページ番号は主な説明箇所を、矢印は別の見出しへの参照を示す。API（Application Programming Interface）は、関数名や引数など、プログラムから機能を呼び出すための取り決めを指す。

\section{実行基盤から性能改善までの章構成}

章構成は、解析作業の順序に対応した四段階からなる。

\begin{enumerate}
\item \textbf{実行基盤の把握}: シェルの章では、現在位置、パス、ファイル、実行ファイル、環境変数、入出力の流れを扱う。続く Python 環境の章の主題は、Python のバージョンとプロジェクトごとの実行環境の管理である。
\item \textbf{基本処理の作成}: Python の基本構文とデータ入出力の章では、短いスクリプトの読解・作成と、ファイル中の値から Python オブジェクトへの変換が中心となる。
\item \textbf{科学データの解析}: NumPy から Astropy までの各章が扱う対象は、数値配列、表、図、統計処理、物理単位、時刻、天体データである。各ライブラリの役割と、処理間のデータ受渡しも比較の対象となる。
\item \textbf{再利用と性能改善}: 自作モジュールと性能改善の章では、解析処理の部品化、テスト、実行時間、メモリ使用量を扱う。総合課題は、ここまでの内容を一つの解析へ統合する位置付けである。
\end{enumerate}

この順序は、特定のライブラリを使うこと自体ではなく、入力を確認してから計算し、結果を検証して保存する流れを身につけるためにある。例えば、pandas で表を読み込めても、ファイルの所在、区切り文字、文字コード、列型を確認できなければ、読み込み結果の妥当性を判断できない。基礎的な操作は、後半の解析を検証するための手段でもある。

\section{コード例と実行結果の読み方}

コード例について、入力、処理、出力、説明との照合を次の順に行う。

\begin{enumerate}
\item \textbf{入力の特定}: 変数、引数、ファイル名、配列の形状、単位
\item \textbf{処理の特定}: 呼び出される関数と、その関数へ渡される引数
\item \textbf{出力の特定}: 返り値、表示内容、保存先のファイル
\item \textbf{説明との照合}: 数値、型、形状、行数、図の軸と、本文が示す想定との一致
\end{enumerate}

例えば、\code{values = np.loadtxt("data.txt")} という行では、\code{data.txt} が入力、\code{np.loadtxt()} が処理、\code{values} が返り値を参照する変数である。\code{values.shape} と \code{values.dtype} の表示から、行列の大きさと要素型を確認できる。各コード例では、入力条件、処理、返り値、実行結果の対応を説明する。

コマンドやコードの実行例には、可能な範囲で代表的な出力を併記した。表示されるパス、バージョン番号、処理時間などは環境によって異なるため、文字列全体ではなく、各行の意味と次の処理へ渡される値を比較する。図については保存済みの出力例を添え、軸、色、凡例、分布、残差の読み方を本文で説明した。

\section{章末課題の進め方}

章末課題では、その章の内容を、それ以前に学んだ操作と組み合わせて使う。前章のファイルを引き継ぐ課題では、指定された入力と作業ディレクトリを確認する。地震カタログの課題は、原則として第~\ref{chap:shell}~章で作る \code{earthquake\_practice} を作業ディレクトリとする。本文中の少数行の CSV は動作確認用の人工データであり、USGS の実際の観測値とは区別する。小さな検証用データで処理を確かめた後、配布データまたは読者自身のデータへ適用する流れが基本となる。

数値の目安は、解答の転記ではなく、コードの誤りを検出するためにある。期待値との不一致があれば、単位、境界条件、並べ替え、欠損値、データ型の順に切り分ける。図を提出する問題では、ファイルの存在だけでは不十分であり、軸ラベルと単位、凡例、対象範囲、データとモデルの対応までが確認項目となる。

第~\ref{chap:kadai}~章の総合課題は、複数章で扱った処理を連結する構成である。課題用スクリプトとデータの入手先は同章に記載した。作業順序は、配布コードによる入出力の把握、小規模データによる正しさの検証、大規模データの処理、性能改善である。

\section{入力・実行条件・結果の検証記録}

データ解析では、コードの正常終了と結果の正しさを区別しなければならない。検証の観点は次の五点である。

\begin{itemize}
\item 手計算可能な小規模入力と、件数、最小値、最大値、平均値などとの比較
\item 配列の \code{shape} と \code{dtype}、表の列名と欠損値数の確認
\item 境界値を含む入力による、\code{<} と \code{<=} などの条件の検証
\item データ、モデル、残差を同じ図に示したフィット結果の確認
\item 高速化前後の結果の一致と、同一条件における処理時間の測定
\end{itemize}

検証に使った入力、コマンド、ライブラリのバージョン、乱数シード、出力ファイルを記録しておけば、後から同じ結果を再現しやすい。解析コードを自作モジュールへ分け、テストを追加する作業も、この記録を実行可能な形にする方法の一つである。

\section{公式資料・付録・参考文献の参照方法}
\label{sec:reference-guide}

本文で取り上げるコマンドやライブラリの使い方は、ソフトウェアの更新によって変わる。本文では概念と代表的な使用例を説明する。指定できる項目の全一覧、関数が返す値、バージョンごとの違いを調べる場合は、公式資料を参照せよ。引用資料は巻末の参考文献にまとめた。

Python の言語仕様と標準ライブラリについては、Python 公式チュートリアル~\cite{python_tutorial} と公式ライブラリリファレンス~\cite{python_library} を参照した。環境構築には JupyterLab、uv、pyenv の公式資料~\cite{jupyterlab_docs,uv_docs,pyenv_docs}、主要な解析ライブラリには各プロジェクトの公式資料~\cite{numpy_quickstart,pandas_tutorials,polars_concepts,matplotlib_quick_start,scipy_user_guide,astropy_user_guide} を用いた。

同じ名前のシェルコマンドでも、macOS と Linux、またはソフトウェアのバージョンによって書き方が異なる場合がある。手元のコンピュータに入っているコマンドの使い方を表示する方法は、第~\ref{chap:shell}~章で具体例とともに説明する。SSH とファイル転送は付録~\ref{app:ssh}、Git と GitHub の基本操作は付録~\ref{app:git} にまとめた。

次章では、データとスクリプトを扱う基盤として、シェルの役割、パス、ファイル操作、入出力、権限を説明する。
